\documentclass[letterpaper, 10 pt, conference]{ieeeconf}
\IEEEoverridecommandlockouts
\overrideIEEEmargins

\usepackage{amsmath,amssymb,amsfonts}
\usepackage{graphicx}
\usepackage{cite}

\title{\bf The GCAT Framework: A Control-Theoretic Model of Governed Autonomy}
\author{[Authors anonymized for submission]}

\begin{document}
\maketitle
\thispagestyle{empty}
\pagestyle{empty}

%%%%%%%%%%%%%%%%%%%%%%%%%%%%%%%%%%%%%%%%%%%%%%%%%%%%%%%%%%%%%%%%%%%%%%%%%%%%%%%%
\begin{abstract}
We introduce the Governed Capability-Autonomy-Trust (GCAT) framework as a control-theoretic model for enforcing admissible autonomy in complex systems. The framework defines a four-dimensional state space capturing governance capacity, control authority, autonomy, and trust. A legitimacy capacity function induces a state-dependent admissible region. We show that this region can be enforced via a control barrier condition and provide a constructive control law. Forward invariance of the admissible set is established under control-affine dynamics with bounded inputs. Boundary behavior and parameter sensitivity are derived explicitly. Simulation results, generated from the stated dynamics, demonstrate that uncontrolled trajectories violate admissibility while barrier-enforced control preserves or restores it.
\end{abstract}

%%%%%%%%%%%%%%%%%%%%%%%%%%%%%%%%%%%%%%%%%%%%%%%%%%%%%%%%%%%%%%%%%%%%%%%%%%%%%%%%
\section{Introduction}

Autonomous systems operate under constraints that extend beyond physical safety into governance, authority, and trust. Existing control frameworks typically enforce physical safety constraints, while governance is treated as external policy or post-hoc analysis.

This work formalizes governance as a state-dependent constraint directly embedded in system dynamics. The central object is an admissible set defined by a capacity-limited inequality linking autonomy to governance, control, and trust. The resulting constraint is enforced at execution through a control barrier condition.

%%%%%%%%%%%%%%%%%%%%%%%%%%%%%%%%%%%%%%%%%%%%%%%%%%%%%%%%%%%%%%%%%%%%%%%%%%%%%%%%
\section{State and Admissibility}

Let the system state be:
\begin{equation}
\mathbf{x} = (g, c, a, t) \in [0,1]^4
\end{equation}

Define the legitimacy capacity:
\begin{equation}
\Lambda(\mathbf{x}) = K g^\alpha c^\beta t^\gamma
\end{equation}

with $K>0$, $\alpha,\beta,\gamma>0$.

Define the admissibility function:
\begin{equation}
h(\mathbf{x}) = \Lambda(\mathbf{x}) - a
\end{equation}

The admissible set is:
\begin{equation}
\mathcal{A} = \{ \mathbf{x} \in [0,1]^4 : h(\mathbf{x}) \ge 0 \}
\end{equation}

%%%%%%%%%%%%%%%%%%%%%%%%%%%%%%%%%%%%%%%%%%%%%%%%%%%%%%%%%%%%%%%%%%%%%%%%%%%%%%%%
\section{Dynamics}

We consider control-affine dynamics:
\begin{equation}
\dot{\mathbf{x}} = f(\mathbf{x}) + G(\mathbf{x}) u
\end{equation}

We instantiate:
\begin{align}
\dot{g} &= k_g (g_0 - g) - \eta_g a \\
\dot{c} &= k_c (c_0 - c) \\
\dot{a} &= k_a a (1-a) + \eta_a (1-g) + u_a \\
\dot{t} &= k_t (t_0 - t)
\end{align}

All parameters are positive constants and $u_a \in \mathbb{R}$ is bounded.

%%%%%%%%%%%%%%%%%%%%%%%%%%%%%%%%%%%%%%%%%%%%%%%%%%%%%%%%%%%%%%%%%%%%%%%%%%%%%%%%
\section{Admissibility Enforcement}

We compute:
\begin{equation}
\dot{h} = \nabla \Lambda(\mathbf{x}) \cdot \dot{\mathbf{x}} - \dot{a}
\end{equation}

where:
\begin{equation}
\nabla \Lambda =
\begin{pmatrix}
\alpha K g^{\alpha-1} c^\beta t^\gamma \\
\beta K g^\alpha c^{\beta-1} t^\gamma \\
\gamma K g^\alpha c^\beta t^{\gamma-1}
\end{pmatrix}
\end{equation}

Substituting dynamics yields:
\begin{align}
\dot{h} &= \alpha K g^{\alpha-1} c^\beta t^\gamma \dot{g}
+ \beta K g^\alpha c^{\beta-1} t^\gamma \dot{c} \nonumber \\
&\quad + \gamma K g^\alpha c^\beta t^{\gamma-1} \dot{t}
- \dot{a}
\end{align}

%%%%%%%%%%%%%%%%%%%%%%%%%%%%%%%%%%%%%%%%%%%%%%%%%%%%%%%%%%%%%%%%%%%%%%%%%%%%%%%%
\section{Forward Invariance}

\textbf{Proposition.} Let $h(\mathbf{x})$ be continuously differentiable and suppose that for all $\mathbf{x} \in \partial \mathcal{A}$ there exists $u_a$ such that:
\begin{equation}
\dot{h}(\mathbf{x}, u_a) \ge 0
\end{equation}

Then $\mathcal{A}$ is forward invariant.

\textit{Argument.} The condition ensures that the vector field does not point outward at the boundary $h=0$. Since the system is locally Lipschitz, trajectories cannot exit $\mathcal{A}$ without violating the boundary condition. Therefore, $\mathcal{A}$ is invariant under admissible control inputs.

%%%%%%%%%%%%%%%%%%%%%%%%%%%%%%%%%%%%%%%%%%%%%%%%%%%%%%%%%%%%%%%%%%%%%%%%%%%%%%%%
\section{Control Law}

We solve for $u_a$ from:
\begin{equation}
\dot{h} \ge -\kappa h
\end{equation}

which yields a constraint:
\begin{equation}
u_a \le \phi(\mathbf{x}) + \kappa h
\end{equation}

where $\phi(\mathbf{x})$ collects all terms independent of $u_a$.

Define:
\begin{equation}
u_a^* = \min(0, \phi(\mathbf{x}) + \kappa h)
\end{equation}

This ensures admissibility is enforced while minimizing intervention.

%%%%%%%%%%%%%%%%%%%%%%%%%%%%%%%%%%%%%%%%%%%%%%%%%%%%%%%%%%%%%%%%%%%%%%%%%%%%%%%%
\section{Boundary Behavior}

At $h=0$:
\begin{equation}
\dot{h} \ge 0
\end{equation}

Substituting yields a constraint linking $\dot{g},\dot{c},\dot{t}$ and $u_a$. This defines the tangent cone of admissible evolution.

%%%%%%%%%%%%%%%%%%%%%%%%%%%%%%%%%%%%%%%%%%%%%%%%%%%%%%%%%%%%%%%%%%%%%%%%%%%%%%%%
\section{Parameter Sensitivity}

\begin{equation}
\frac{\partial \Lambda}{\partial g} = \alpha K g^{\alpha-1} c^\beta t^\gamma
\end{equation}

Sensitivity increases with multiplicative capacity, implying that small changes in governance or trust can have amplified effects on admissibility near high-capacity regimes.

%%%%%%%%%%%%%%%%%%%%%%%%%%%%%%%%%%%%%%%%%%%%%%%%%%%%%%%%%%%%%%%%%%%%%%%%%%%%%%%%
\section{Simulation}

The system is simulated using the defined dynamics. Initial conditions are chosen such that $h(\mathbf{x}_0) > 0$ but near the boundary.

Uncontrolled trajectories ($u_a=0$) result in $h(\mathbf{x}(t)) < 0$ within finite time.

Controlled trajectories ($u_a = u_a^*$) maintain $h(\mathbf{x}(t)) \ge 0$ for all $t$.

\begin{figure}[t]
\centering
\includegraphics[width=\linewidth]{gcat_phase_boundary_trajectories.pdf}
\caption{Phase boundary and trajectories derived from the stated dynamics.}
\label{fig:phase}
\end{figure}

%%%%%%%%%%%%%%%%%%%%%%%%%%%%%%%%%%%%%%%%%%%%%%%%%%%%%%%%%%%%%%%%%%%%%%%%%%%%%%%%
\section{Related Work}

Control barrier functions provide a framework for enforcing state constraints in nonlinear systems. This work extends that framework to constraints derived from governance and system capacity.

%%%%%%%%%%%%%%%%%%%%%%%%%%%%%%%%%%%%%%%%%%%%%%%%%%%%%%%%%%%%%%%%%%%%%%%%%%%%%%%%
\section{Conclusion}

We presented a control-theoretic formulation of governed autonomy. The admissible region is enforced through a barrier condition, ensuring that autonomy remains bounded by system capacity.

\end{document}
